%----------------------------------------------------------------------
% Appendix: decomposing "Syrian insurgents" with a random forest
%----------------------------------------------------------------------
\subsection{Decomposing the label \emph{Syrian insurgents}}
\label{app:rf}

$63{,}791$ events in the merged data carry the actor label \emph{Syrian
insurgents}, which in UCDP's coding stands for one or more opposition
formations that the reporting does not identify individually. We assign these
events to the three formations that account for most of the named opposition
events, Ahrar al-Sham, HTS (Hayat Tahrir al-Sham) and the Free Syrian Army,
using a random forest trained on the events where the actor is named.

\paragraph{Training sample and features.}
The training sample consists of the $9{,}830$ events in which at least one of the
three formations appears in one of the actor columns \texttt{Side\_a\_*} or
\texttt{Side\_b\_*}. Since several of them can fight in the same event, the
target is not a single class but three binary indicators, one per formation.
The predictors are year, calendar month, latitude and longitude. We use no
further event characteristics: source, violence type and fatality counts are
coded differently in UCDP and ACLED, so the model would partly learn the source
rather than the actor.

\paragraph{Model and performance.}
We fit a random forest with $10$ trees on $70\%$ of the labelled events
($6{,}881$) and evaluate it on the remaining $30\%$ ($2{,}949$); the random seed
is fixed. Table~\ref{tab:rf_performance} reports precision, recall, $F_1$ and
accuracy for the positive class, that is, for correctly identifying that a
formation was present. Accuracy lies between $0.90$ and $0.93$, with precision
and recall of similar magnitude for all three formations.

\begin{table}[htbp]
\centering
\caption{Out-of-sample performance of the random forest (positive class)}
\label{tab:rf_performance}
\begin{tabular}{lcccc}
\hline\hline
Formation & Precision & Recall & $F_1$ & Accuracy \\
\hline
Ahrar al-Sham             & $0.89$ & $0.84$ & $0.86$ & $0.92$ \\
HTS: Hayat Tahrir al-Sham & $0.94$ & $0.94$ & $0.94$ & $0.93$ \\
Free Syrian Army          & $0.89$ & $0.86$ & $0.87$ & $0.90$ \\
\hline\hline
\end{tabular}
\end{table}

\paragraph{Assignment rule.}
For each unnamed event $i$ the forest returns a probability $\hat{p}_{ig}$ that
formation $g$ took part in it. Formation $g$ is assigned to event $i$ whenever
$\hat{p}_{ig}\geq\tau$, with $\tau=0.7$. Events for which no formation reaches
the threshold are left unassigned: we prefer to lose them rather than to impose
an actor the model is not confident about. This applies to $7{,}490$ of the
$63{,}791$ events ($11.7\%$). Of the $56{,}301$ events that are assigned,
$49{,}122$ receive one formation, $6{,}513$ two and $666$ all three. Expanding
the multiple assignments into one observation per formation yields $64{,}146$
actor level observations.

\paragraph{Plausibility check.}
The true actor is unobserved, so the predictions are compared with the area in
which each formation operated in a given year. For every formation and year we
keep the named events with at least $k=3$ other such events within $l=3$~km and
take the convex hull of the retained points. $20{,}682$ of the $56{,}301$ assigned
events fall inside at least one of these polygons; for $17{,}544$ of them
($84.8\%$) the predicted formation is one whose polygon covers the event
(Table~\ref{tab:rf_polygon}). Since the polygons are constructed from the same
named events used for training, this is a consistency check rather than an
independent validation.

\begin{table}[htbp]
\centering
\caption{Predicted formations and their yearly areas of operation}
\label{tab:rf_polygon}
\begin{tabular}{lccc}
\hline\hline
Predicted formation & Predictions & Inside polygon & Share \\
\hline
Ahrar al-Sham             & $8{,}884$  & $8{,}649$ & $0.97$ \\
HTS: Hayat Tahrir al-Sham & $13{,}033$ & $8{,}716$ & $0.67$ \\
Free Syrian Army          & $2{,}324$  & $1{,}081$ & $0.47$ \\
\hline\hline
\end{tabular}
\end{table}
